PHQD denoises Gaussian noise in the latent space of an image rather than in pixel space, and it does so one small patch at a time so that each patch fits a few-qubit circuit. Figure~\ref{fig:architecture} shows the workflow. A Fashion-MNIST image is compressed by the encoder of a convolutional autoencoder into a latent map, the map is corrupted with Gaussian noise, and the noisy map is cut into overlapping patches. Each patch is denoised by a hybrid quantum--classical model, the denoised patches are folded back into a latent map by averaging, an optional stitching network refines the result, and the decoder produces the output image. The pipeline is implemented as five modules: the autoencoder, the quantum patch encoder, the stitching network, the classical baselines, and the quantum denoisers with the comparison routine.

We use Fashion-MNIST~\cite{fmnist}, which contains 60,000 training and 10,000 test grayscale images of $28\times28$ pixels in ten clothing classes. Pixel values are converted to tensors and normalised with mean 0.5 and standard deviation 0.5, so that they lie in $[-1,1]$, which matches the Tanh output of the decoder. The autoencoder is trained first with an MSE reconstruction loss, Adam at learning rate $10^{-3}$, cosine annealing, batch size 128 and 20 epochs, and is then frozen. Its encoder acts as the feature extractor: it maps every image to a latent map $z\in\mathbb{R}^{64\times7\times7}$ with 3,136 values. For denoising, the latent map is corrupted as $\tilde z=z+\sigma\epsilon$ with $\epsilon\sim\mathcal{N}(0,I)$ and $\sigma=0.3$, and a $2\times2$ window with stride one is slid over the $7\times7$ grid to give 36 overlapping patches of 256 values each.

The proposed denoiser maps a noisy patch $p$ to an estimate $\hat p$ in three steps. A classical pre-projection first reduces $p$ to $n$ values that are bounded by a Tanh function and rescaled to rotation angles in $[0,\pi]$. A PQC then encodes these angles with $R_Y$ rotations, applies layers of cyclic CNOT gates and trainable rotations, and returns the Pauli-$Z$ expectation value of each qubit, giving the vector $q(p)$. A classical post-projection maps $q(p)$ back to patch space, and a learned linear residual path from the input is added, so that $\hat p=g_{\mathrm{post}}(q(p))+W_rp+b_r$. Four 4-qubit configurations differ in circuit depth (1, 2 or 3 layers) and in data encoding (angle encoding or an IQP-style encoding), as listed in Table~\ref{tab:configs}. A 6-qubit variant, EPQD-6Q-V2, re-uploads the input in every layer, reads out the $X$, $Y$ and $Z$ expectation values of all qubits, which gives 18 circuit features, and is trained jointly with the stitching network. The circuits are simulated with PennyLane on a classical state-vector simulator, and gradients are obtained by back-propagation through the simulator.

Training follows a classical warm-up and quantum fine-tuning protocol. In the standard models the quantum weights are first frozen and a linear bypass replaces the circuit, so that the classical projections are trained for 10 epochs at learning rate $10^{-3}$. The bypass is then removed and the whole model is fine-tuned for 8 epochs at learning rate $3\times10^{-3}$ on at most 8,000 training patches. Both stages use the MSE loss between the denoised and clean patch, Adam, cosine annealing and batch size 128. EPQD-6Q-V2 is trained in three phases of 15, 8 and 6 epochs with learning rates $10^{-3}$, $2\times10^{-3}$ and $5\times10^{-4}$, and the last phase minimises a combination of latent MSE, image MSE with weight 0.5 and a structural similarity term, $1-\mathrm{SSIM}$, with weight 0.2.

The hybrid models are compared with a patch-wise MLP and a U-Net that works on the whole latent map. Both are trained with the same noise level, loss and optimiser settings for 20 epochs, and the autoencoder reconstruction of the clean latent map and the noisy input serve as the upper and lower reference. All methods see the same noisy latent maps, so differences in the results come from the denoiser and not from the data. Quality is measured on the Fashion-MNIST test split by the MSE, PSNR and SSIM between the decoded and original images, and by the Fr\'echet Inception Distance computed with Inception-V3 features on 1,000 test images. The remaining subsections describe each module, its architecture and its parameters in detail.
